% Created 2018-10-01 Mon 15:38
% Intended LaTeX compiler: pdflatex
\documentclass[11pt]{article}
\usepackage[utf8]{inputenc}
\usepackage[T1]{fontenc}
\usepackage{graphicx}
\usepackage{grffile}
\usepackage{longtable}
\usepackage{wrapfig}
\usepackage{rotating}
\usepackage[normalem]{ulem}
\usepackage{amsmath}
\usepackage{textcomp}
\usepackage{amssymb}
\usepackage{capt-of}
\usepackage{hyperref}
\usepackage[margin=0.80in]{geometry}
\usepackage[backend=biber, style=ieee, url=false]{biblatex}
\usepackage{float}
\usepackage[super,negative]{nth}
\usepackage[capitalise]{cleveref}
\usepackage{pst-node,transparent,ragged2e}
\addbibresource{/home/nlisa/.spacemacs.d/org-files/bibliography.bib}
\DeclareFieldFormat[inproceedings]{citetitle}{\textit{#1}}
\DeclareFieldFormat[inproceedings]{title}{\textit{#1}}
\DeclareFieldFormat[misc]{citetitle}{#1}
\DeclareFieldFormat[misc]{title}{#1}
\renewcommand*{\bibpagespunct}{%
\ifentrytype{inproceedings}
{\addspace}
{\addcomma\space}}
\AtEveryCitekey{\ifuseauthor{}{\clearname{author}}}
\AtEveryBibitem{\ifuseauthor{}{\clearname{author}}}
\SpecialCoor
\author{Nyameko Lisa}
\date{30 September 2018}
\title{IPSP089 - Legal Aspects of Traditional Knowledge and Biodiversity\\\medskip
\large Assignment 2: 656029}
\hypersetup{
 pdfauthor={Nyameko Lisa},
 pdftitle={IPSP089 - Legal Aspects of Traditional Knowledge and Biodiversity},
 pdfkeywords={},
 pdfsubject={},
 pdfcreator={Emacs 26.1 (Org mode 9.1.14)},
 pdflang={English}}
\begin{document}

\maketitle
\addvspace{110pt}
\centering{
\pnode(0.5\textwidth,-0.5\textheight){thisCenter}
\rput(thisCenter){%\transparent{0.25}
\includegraphics[width=2.7in]{/home/nlisa/course/llb/wipo-unisa/UNISACoatofArms.eps}}}

\justifying
\addvspace{110pt}
\section*{Declaration}
\label{sec:orgbd182c5}
I know that plagiarism is to use someone else’s work and pass it off as my own.\\
I know that plagiarism is wrong.\\
I confirm that this assignment is my own work.\\
I have acknowledged in the bibliography accompanying the assignment all the sources that I have used.\\
I have not directly copied without acknowledgement anything from the Internet or from any other source.\\
I have indicated every quotation and citation in a footnote or bracket linked to that quotation.\\
I have not allowed anyone else to copy my work and to pass it off as their own work.\\
I understand that if any unacknowledged copying whatsoever appears in my assignment I will receive zero per cent for the assignment.\\
I am aware of the UNISA policy on plagiarism and understand that disciplinary proceedings can be instituted against me by UNISA if I contravene this policy.\\
I indicate my understanding and acceptance of this declaration by
entering my name here under:
\begin{itemize}
\item Name: \textbf{Nyameko Lisa} (Student Number: \textbf{7874-909-3})
\end{itemize}

\subsection*{NOTE}
\label{sec:org1f0d708}
Please note that footnotes will be denoted as \footnote{This is a footnote.} and will
appear at the bottom of the page.\\
References will be denoted by \cite{rsa93_tm_act} and will appear at the end of the document.
\newpage

\section{Question 1 [50]}
\label{sec:orgd8bad80}
Discuss the following in respect of traditional cultural expressions
(expressions of folklore):

\subsection{Why is there a need to protect expressions of folklore? (20)}
\label{sec:org4f61271}

Folklore is an important element of a nation's cultural heritage, providing its
people with an avenue for self-expression and it exists as a means of a
country's social identity. Traditional expressions have generally been
characterized as a broad set of oral and or written literary, artistic,
religious, scientific, technological works and other traditions and productions
that are transmitted from one generation to another, and whose creators need not
necessarily be indigenous peoples.\\

With regards to folklore, aspects related to its subject matter, technological
and methodological content need to be taken into consideration. Thus it is
important to consider not only protection from a copyright perspective, but also
the mechanisms of industrial property protection, including patents, trademarks
and industrial designs.\\

The need for protection of expressions of folklore arises where the creators or
rights holders, would need to prevent others from reproducing their creations or
disseminating them to the public through performance or broadcast, \emph{(i.e. law of
copyright)}; and to restrict the unauthorized misappropriation of technological
and methodological ideas, inventions, crafts, designs and carvings, \emph{(i.e. laws
of patents, trademarks and industrial designs)}. Moreover, informal traditional
customary regimes may not provide adequate protection against cultural, economic
or commercial exploitation.\\

Both globalization and the advent of technological telecommunications advances
in recent history, have lead to the improper exploitation of cultural
expressions of folklore currently experienced today. The misappropriation of the
intellectual property rights that \emph{`should'} be enjoyed by a people, resulting
from their endeavors and accumulated cultural heritage can take a number of
forms, from commercial exploitation to adaptation and distortion of the original
works.\\

There exists a clear need for the development of both national and international
instruments $\backslash$ frameworks, providing legal protection against the exploitation of
cultural expressions of folklore, without any compensation to their originating
communities. Moreover, litigation processes are also required to reclassify
those expressions of folklore that were previously considered to belong to the
public domain. Not forgetting the subsequent infringement and sanctioning
procedures with regards to infringement of the rights in `works of folklore'.

\subsection{The WIPO / UNESCO's Model Provisions on the issue. (30)}
\label{sec:org7d08ce3}

As per the provisions of [article 15(4)(a-b)]\cite{wipo86_berne}, although not
explicitly limited to, nor directly referring to folklore, member Union
countries may use their respective national legislation to designate a competent
authority to represent, protect and enforce the rights of an unknown, but
assumed to be a domiciled in said country, author of an unpublished work.
After which, that member of the Union must by way of written declaration notify
the WIPO Director General of said details, who in turn will communicate the
details of the declaration to all other countries of the Union, as per the
introductory observations of \cite{wipo85_model_provi_national_laws_folklore}.\\

In 1980 a working group was established between WIPO and UNESCO, where a study
for the draft model provisions for national laws on the protection of creations
of folklore were undertaken. In 1982 the Directors General of UNESCO and WIPO
convened a Committee of Governmental Experts on the Intellectual Property
Aspects of the Protection of Expressions of Folklore, where the \emph{`Model Provisions for National Laws on the Protection of Expressions of Folklore
Against Illicit Exploitation and Other Prejudicial Actions}' were adopted.
Moreover these provisions were designed in such a way as to be adaptable to
appropriate national conditions, with the primary objective of protecting
against the abuse [section 1]\cite{wipo85_model_provi_national_laws_folklore}, of
expressions of folklore whilst also encouraging it's further development and
dissemination, as per the introductory observations of
\cite{wipo85_model_provi_national_laws_folklore}.\\

[Section 2]\cite{wipo85_model_provi_national_laws_folklore} articulates the
expressions of folklore protected under these Model Provisions, as
\emph{``productions consisting of characteristic elements of the traditional artistic
heritage developed and maintain by a community''}, including verbal, musical,
physically demonstrative and tangible expressions. [Sections 3 and
4]\cite{wipo85_model_provi_national_laws_folklore} deal with the authorized
utilization and corresponding exceptions respectively. Permission from a
competent authority is required for: any publication, reproduction or any
distribution of copies of [section
3(i)]\cite{wipo85_model_provi_national_laws_folklore}; or any public recitation,
performance or transmission [section
3(ii)]\cite{wipo85_model_provi_national_laws_folklore} of expressions of folklore.
With the exception of utilization that is educational [section
4.1(i)]\cite{wipo85_model_provi_national_laws_folklore}; an illustration
consistent with fair practice [section
4.1(ii)]\cite{wipo85_model_provi_national_laws_folklore}; borrowed and re-used in
the creation of an original work of an author(s) [section
4.1(iii)]\cite{wipo85_model_provi_national_laws_folklore}; or where the
utilization is incidental [section
4.2]\cite{wipo85_model_provi_national_laws_folklore}.\\

Utilization of an expression of folklore without either consent from a competent
authority nor appropriate acknowledgement of source of the author or community
and / or geographic place from which the work originates [section
5.1]\cite{wipo85_model_provi_national_laws_folklore}, constitute offensive
violations [section 6]\cite{wipo85_model_provi_national_laws_folklore} which are
liable to seizure or other actions [section
7]\cite{wipo85_model_provi_national_laws_folklore} and other civil remedies such
as damages [section 8]\cite{wipo85_model_provi_national_laws_folklore}. Where the
definition of a `competent or supervisory authority' [section
9]\cite{wipo85_model_provi_national_laws_folklore} is left open to the relevant
member Union country.\\

[Section 10[\cite{wipo85_model_provi_national_laws_folklore} details the necessary
procedures required for authorization with respect to utilization of an
expression of folklore. Such requests should be in writing to the competent
authority / community concerned [section
10.1]\cite{wipo85_model_provi_national_laws_folklore}. Where a fee may be
collected, as established by the supervisory authority, for the purposes of
promoting / safeguarding national culture and folklore [section
10.2]\cite{wipo85_model_provi_national_laws_folklore}. The Model Provisions also
accommodate appeals from authorization applicant against decisions made by the
competent authority [section
10.3]\cite{wipo85_model_provi_national_laws_folklore}.\\

As per the provisions of [section
11]\cite{wipo85_model_provi_national_laws_folklore}, with respect to matters
relating to appeals or offenses, the question of jurisdiction is left open and
to be defined by the relevant member Union country. This vague and opened-ended
definition in the Model Provisions is indeed problematic and inadequate,
especially in light of electronic transactions and the nature of potential
infringements and offenses that span multiple jurisdictions, across a number of
member Union countries. As per the provisions of [section
12]\cite{wipo85_model_provi_national_laws_folklore} the Model Provisions shall in
no way limit, prejudice nor supersede any protections applicable to expressions
of folklore, enjoyed under copyright law; industrial property law; laws
protecting producers, performers and broadcasting organizations; or any other
laws or international treaties to which the country is a party to. Lastly
[sections 13 and 14]\cite{wipo85_model_provi_national_laws_folklore} describe how
the Model Provisions cannot be interpreted in any way to hinder the use or
development of expressions of folklore, and where expressions of folklore are
developed in a foreign country they are subject to reciprocal or international
treaty based protection.


\printbibliography
\end{document}